\documentclass[11pt]{article}
\usepackage[T1]{fontenc}
\usepackage{lmodern,microtype,amsmath,amssymb,amsthm,booktabs}
\usepackage[margin=1in]{geometry}
\usepackage[authoryear,round]{natbib}
\usepackage[hidelinks]{hyperref}
\hypersetup{pdftitle={Sharp subsidies for protecting core maxima in three-sided assignment markets},pdfauthor={Author identity pending confirmation}}
\newtheorem{theorem}{Theorem}[section]
\newtheorem{lemma}[theorem]{Lemma}
\newtheorem{proposition}[theorem]{Proposition}
\newtheorem{corollary}[theorem]{Corollary}
\theoremstyle{remark}
\newtheorem{remark}[theorem]{Remark}
\newcommand{\R}{\mathbb R}
\newcommand{\core}{\mathcal C}
\newcommand{\pack}{\mathcal P}
\newcommand{\one}{\mathbf 1}
\newcommand{\E}{\mathbb E}
\DeclareMathOperator{\supp}{supp}
\title{Sharp subsidies for protecting core maxima\\in three-sided assignment markets}
\author{Author name(s): confirmation required}
\date{17 October 2026\\[3pt]\small International Journal of Game Theory candidate; not submitted}
\begin{document}
\maketitle

\begin{abstract}
Consider three equally sized finite sectors, nonnegative pair values, and
productive teams containing exactly one player from each sector. A team's
value is the sum of its three pair values; coalition worth is maximum integral
matching worth. Assume the original core is nonempty and grand worth is
$T>0$. For each player let $H_i$ be their maximum payoff over that entire
original core. We determine the worst-case external subsidy needed to cover
every original coalition while paying each player at least $H_i/3$.
Every finite market requires strictly less than $T/3$, and the supremum of
the subsidy-to-worth ratio is exactly $1/3$. A direct fractional-packing
argument combines six core-maximum comparisons with residual budgets at
separately attaining core allocations. For an integral partial matching
leaving $k\ge1$ players unmatched in each sector, the associated subsidy
obstruction is at most $2(k-1)T/[3(2k-1)]$. This bound is sharp for every
$k\ge2$. We give complete proofs, a family with pair values in
$\{0,1,2\}$, and reproducible exact checks. A compact supplement records exact nonnegative slack identities and the
verification scope. The original core is fixed before subsidy, and the result
is static: it makes no claim about reporting incentives across profiles.
\end{abstract}
\noindent\textbf{Keywords:} assignment game; core; individual payoff maximum;
subsidy; cost of stability; fractional matching.

\section{Introduction}
A stable allocation need not simultaneously give each player a substantial
fraction of what that player could obtain at another stable allocation.
An external budget permits a different question: how much must be added to
the grand coalition's worth to protect individual benchmarks while retaining
all original coalition constraints?

We study the benchmark $H_i/3$, where $H_i$ is player $i$'s maximum over the
entire original core of a three-sided assignment market. The reference core
is fixed before intervention. In particular, the benchmark is not recomputed
after adding the subsidy. The restriction to nonnegative pair-additive team
values is central; an arbitrary three-dimensional value array is not assumed
to satisfy our conclusion.

Our main result gives the sharp worst-case subsidy ratio $1/3$. It is a
supremum over finite markets of arbitrary size, not a maximum at a finite
size. The upper proof applies directly to fractional partial matchings in
the subsidy dual. It does not require a reduction to sparse markets,
replication of players, or an acyclic support. A second result determines a
sharp upper bound for the obstruction supplied by an integral partial
matching, indexed by its number of unmatched players per sector.

The proof uses two estimates for the same dual expression. If $W$ is a
fractional matching's value, $S$ the sum of residual-weighted individual
maxima and $\beta$ a residual-weighted diagonal pair budget, then
\[
                 3W+S\le5T-\beta,
        \qquad   3W+S\le3T+\beta.
\]
The first estimate aggregates six comparisons between individual maxima;
the second uses a different original-core witness for each maximum.
Averaging eliminates $\beta$. Retaining the slack terms proves strictness in
every finite market and the sharper integral bound.

\paragraph{Scope of this article.}
The main paper contains the full three-sided theorem, strict finite-market
upper bound, integral obstruction refinement and matching family. Online
Resource~1 records exact slack identities and reproducibility details.
The project's broader October~7 compendium also contains higher-sector,
flow, transportation and ancillary core results. Those separate extensions
are not included or claimed as additional contributions of this article.

\paragraph{Related work and contribution.}
The assignment-game connection between core allocations and linear-programming
duality is classical~\citep{ShapleyShubik1971}. Protection relative to individual
core maxima is also an established question. N\'u\~nez and Rafels (2004, \emph{Bargained stable allocations in assignment
markets}, Barcelona Economics Working Paper~153, Definition~6;
\href{https://www.bse.eu/sites/default/files/working_paper_pdfs/153.pdf}{inspected working version}) study the bilateral half-optimal core:
stable allocations paying every agent at least half their maximum core
payoff. Thus neither the benchmark based on core maxima nor simultaneous
fractional protection is introduced here.

External payments to stabilize a game are studied through the cost of
stability (Bachrach et al., 2009, arXiv:0907.4385; published successor
\citep{Bachrach2018}). Meir, Rosenschein and Malizia
\citep[Section~2.2]{Meir2011} also define the cost of extending a given
efficient payoff vector to a stable one; the same definition appears in
Bachrach et al.~\citep[Section~5]{Bachrach2018}. Their input vector has total $T$;
our floor vector $H/3$ may have smaller total. Section~\ref{sec:framework}
records the exact comparison. Priority-game transformations provide another
precedent for adding individual claims~\citep{Priority2024}.

Atay and N\'u\~nez~\citep[p.~273]{AtayNunez2019} describe a model in
Quint~\citep{Quint1991} in which full teams have pair-additive values and
smaller coalitions have zero worth. Their own graph model permits productive
partial-sector teams and therefore has different coalition constraints.
Results on core extrema in supermodular multisided markets with local
monotonicity~\citep{Martinez2026} impose additional hypotheses absent here.

The contribution proposed for review is the aggregate compatibility
inequality and the resulting sharp subsidy bound for the stated class.
The extremal family was obtained earlier in this project's subsidy work;
the matching upper bound and its integral refinement complete its extremal
interpretation. This manuscript gives a self-contained proof rather than
claiming a new general subsidy framework or a new assignment model.
A companion manuscript, \emph{Acyclicity and simultaneous protection of
individual core maxima}, asks which fractions of the same benchmarks can
be protected \emph{within} the original core and studies graph structure.
Its triangle construction differs from the uniform-ideal family used here.
The present paper fixes the fraction at one third and optimizes external
subsidy. Another companion, \emph{Sharp subsidies for truthful production
and network cooperation}, links payments across report profiles; this
static result supplies no such incentive guarantee.
Historical priority of the precise result remains under review. In
particular, full-original comparisons with Quint~\citep{Quint1991},
Lucas~\citep{Lucas1995} and Stuart~\citep{Stuart1997} remain incomplete.
Other original-source leads and general implication routes are listed in
the companion source ledger; no exhaustive priority clearance is claimed.

\section{Model and subsidy duality}\label{sec:model}
Let $N_1,N_2,N_3$ be disjoint sectors of size $n\ge1$, and put
$N=N_1\cup N_2\cup N_3$. A market $M$ is specified by nonnegative real
matrices $A,B,C$ on $N_1\times N_2$, $N_1\times N_3$, and
$N_2\times N_3$, respectively. A full team $e=(i_1,i_2,i_3)$ has value
\begin{equation}\label{eq:team}
                  v(e)=A_{i_1i_2}+B_{i_1i_3}+C_{i_2i_3}.
\end{equation}
For a coalition $Q\subseteq N$, its worth $v(Q)$ is the maximum sum of
values of player-disjoint full teams contained in $Q$. The empty matching
is allowed. Pairs and all other coalitions containing no full team have
worth zero, even when a constituent pair matrix entry is positive.

Write $T=v(N)>0$. We assume throughout that the original core
\[
 \core=\{x\in\R^N:x(N)=T,\ x(Q)\ge v(Q)\text{ for every }Q\subseteq N\}
\]
is nonempty, where $x(Q)=\sum_{i\in Q}x_i$.
Nonnegativity and full-team coverage describe all coalition constraints:
\begin{equation}\label{eq:core}
 \core=\{x\in\R_+^N:x(N)=T,\ x(e)\ge v(e)\text{ for every full team }e\}.
\end{equation}
Indeed, sum the full-team inequalities over an optimal integral matching
inside any coalition and use nonnegative payoffs for unmatched players.
The reverse implication follows from singleton and full-team coalitions.
The core is a compact subset of the nonnegative total-$T$ simplex. Define
the attained individual maxima
\begin{equation}\label{eq:ideals}
                        H_i=\max_{x\in\core}x_i.
\end{equation}

The protected subsidy is
\begin{equation}\label{eq:subsidy}
 s(M)=\min\{y(N)-T:y_i\ge H_i/3\ (i\in N),\
                         y(e)\ge v(e)\ (e\text{ a full team})\}.
\end{equation}
This feasible linear program has a finite minimum. Each feasible vector
is nonnegative and covers all original coalitions, including $N$, so
$s(M)\ge0$. It is interpreted as an allocation for a subsidized grand
coalition; it need not belong to the original core.

Let $\pack$ be the fractional partial-matching polytope:
\[
 \pack=\left\{\lambda\ge0:
               p_i:=\sum_{e\ni i}\lambda_e\le1\quad(i\in N)\right\}.
\]
For $\lambda\in\pack$ put $q_i=1-p_i$ and define
\begin{equation}\label{eq:dualquantities}
 W=\sum_e\lambda_e v(e),\qquad
 S=\sum_iq_iH_i,\qquad V=3W+S.
\end{equation}

\begin{lemma}[Protected subsidy dual]\label{lem:dual}
For every market satisfying the preceding assumptions,
\begin{equation}\label{eq:dual}
                  s(M)=\max_{\lambda\in\pack}(W+S/3-T)
                       =\max_{\lambda\in\pack}(V/3-T).
\end{equation}
Both extrema are attained.
\end{lemma}
\begin{proof}
Substitute $y_i=H_i/3+z_i$ in~\eqref{eq:subsidy}. The variable part is
to minimize $\sum_i z_i$ subject to $z_i\ge0$ and
$\sum_{i\in e}z_i\ge v(e)-\sum_{i\in e}H_i/3$ for every team $e$.
Its dual has one nonnegative variable $\lambda_e$ per team and constraints
$\sum_{e\ni i}\lambda_e\le1$. Adding the constant
$\sum_iH_i/3-T$ to the dual objective gives~\eqref{eq:dual}.
The primal is feasible and bounded below. Moreover, $\pack$ is nonempty
and compact, since there are finitely many teams and $0\le\lambda_e\le1$.
Finite-dimensional linear-programming duality proves the assertion, for
real as well as rational market data.
\end{proof}

\begin{theorem}[Sharp subsidy ratio]\label{thm:main}
Every market in the stated class satisfies $s(M)<T/3$. Nevertheless,
\begin{equation}\label{eq:sharp}
             S_3:=\sup_M\frac{s(M)}{v(N)}=\frac13.
\end{equation}
The supremum is approached by markets with pair values in $\{0,1,2\}$.
\end{theorem}
We prove the upper bound in Section~\ref{sec:aggregation} and construct the
sharpness family in Section~\ref{sec:family}.

\section{Diagonal budgets and comparisons of maxima}\label{sec:comparisons}
Choose an optimal integral matching. Because sector sizes are equal, it
can be completed to a perfect matching; nonnegative team values imply
that the completion remains optimal. Relabel the players so that this
perfect matching consists of $(1,j),(2,j),(3,j)$ for $j=0,\ldots,n-1$.
Let
\begin{align}
 d_j&=A_{jj}+B_{jj}+C_{jj}, &\sum_jd_j&=T,\label{eq:diagonal}\\
 u_{1,j}&=A_{jj}+B_{jj},&u_{2,j}&=A_{jj}+C_{jj},
 &u_{3,j}&=B_{jj}+C_{jj}.\notag
\end{align}
Thus $u_{a,j}$ is the diagonal pair budget incident to player $(a,j)$, and
\begin{equation}\label{eq:budget}
                     \sum_{a,j}u_{a,j}=2T.
\end{equation}
Every diagonal team is tight at every $x\in\core$: their covering
inequalities sum to the efficiency equality. Consequently
\begin{equation}\label{eq:caps}
     \sum_{a=1}^3x_{a,j}=d_j,\qquad
     0\le H_{a,j}\le d_j,\qquad
     \delta_{a,j}:=d_j-H_{a,j}\ge0.
\end{equation}
Zero diagonal values cause no difficulty; no argument divides by $d_j$.

\begin{lemma}[Same-sector range]\label{lem:range}
For every $x\in\core$, sector $a$, and labels $j,h$,
\begin{equation}\label{eq:range}
                         x_{a,j}-x_{a,h}\le u_{a,j}.
\end{equation}
\end{lemma}
\begin{proof}
Complete the two players of the diagonal team at $j$ outside sector $a$
with player $(a,h)$. By pair nonnegativity this team's value is at least
the opposite diagonal pair value, namely $d_j-u_{a,j}$.
Core coverage and diagonal tightness give
$x_{a,h}+d_j-x_{a,j}\ge d_j-u_{a,j}$.
This is~\eqref{eq:range}; it also holds when $j=h$.
\end{proof}

\begin{lemma}[Six ideal comparisons]\label{lem:comparison}
For a full team $e=((1,j_1),(2,j_2),(3,j_3))$ and every ordered
permutation $(a,b,c)$ of $(1,2,3)$,
\begin{equation}\label{eq:comparison}
         H_{c,j_c}-H_{c,j_b}\ge v(e)-d_{j_b}-u_{a,j_a}.
\end{equation}
The labels $j_1,j_2,j_3$ need not be distinct.
\end{lemma}
\begin{proof}
Subtract diagonal tightness at label $j_b$ from coverage of $e$.
For every original core point $x$ this gives
\[
 x_{c,j_c}-x_{c,j_b}
 \ge v(e)-d_{j_b}+x_{a,j_b}-x_{a,j_a}
 \ge v(e)-d_{j_b}-u_{a,j_a},
\]
using Lemma~\ref{lem:range}. Choose a core point attaining the
\emph{subtracted} maximum $H_{c,j_b}$. Its other coordinate is at most
$H_{c,j_c}$, proving~\eqref{eq:comparison}. A different core point may
be used for every comparison. No joint attainment is assumed.
\end{proof}

\section{Aggregation for arbitrary fractional matchings}\label{sec:aggregation}
Fix $\lambda\in\pack$. In addition to~\eqref{eq:dualquantities}, write
\begin{equation}\label{eq:residual}
 k=n-\sum_e\lambda_e,\qquad
 \beta=\sum_iq_iu_i,\qquad C_2=\sum_iq_i^2u_i.
\end{equation}
Here and below $u_i=u_{a,j}$ when $i=(a,j)$. Every sector has the same
residual mass: $\sum_{i\in N_a}q_i=k\ge0$. In particular $\sum_iq_i=3k$.
For each diagonal label let
\[
 n_j=\sum_{a=1}^3p_{a,j},\qquad H_j^\Sigma=\sum_{a=1}^3H_{a,j},\qquad
 E=\sum_{a,j}\left(\sum_{b\ne a}q_{b,j}\right)\delta_{a,j}\ge0.
\]
The symbol $n_j$ is a fractional load on a diagonal label, not the number
of labels $n$; it may be any number in $[0,3]$.

\begin{lemma}[Aggregate ideal comparison]\label{lem:aggregate}
Every fractional partial matching satisfies
\begin{equation}\label{eq:aggregate}
                          V\le5T-\beta-\frac E2.
\end{equation}
\end{lemma}
\begin{proof}
Sum~\eqref{eq:comparison} over the six ordered sector permutations for
one team $e$. Each participating ideal appears twice positively. The
negative cross terms are
$-\sum_{a=1}^3 H_{j_a}^\Sigma+\sum_{a=1}^3H_{a,j_a}$.
Each diagonal value and each incident budget occurs twice. Therefore
\[
 6v(e)\le2\sum_ad_{j_a}+2\sum_au_{a,j_a}
             +3\sum_aH_{a,j_a}-\sum_aH_{j_a}^\Sigma.
\]
Multiply by $\lambda_e$ and sum over teams. With
$D_p=\sum_{a,j}p_{a,j}d_j$, $U_p=\sum_{a,j}p_{a,j}u_{a,j}$ and
$H_p=\sum_{a,j}p_{a,j}H_{a,j}$, this becomes
\begin{equation}\label{eq:sixsum}
                     6W\le2D_p+2U_p+3H_p-\sum_jn_jH_j^\Sigma.
\end{equation}
All sums retain sector multiplicities, including when team labels coincide.

Add $2S$ to~\eqref{eq:sixsum}. For one label $j$, the coefficient of
$H_{a,j}$ is
\begin{equation}\label{eq:coefficient}
 3p_{a,j}+2q_{a,j}-n_j
       =2+p_{a,j}-n_j=\sum_{b\ne a}q_{b,j}.
\end{equation}
Replace $H_{a,j}$ by $d_j-\delta_{a,j}$. Summing
\eqref{eq:coefficient} over $a$ gives $2\sum_aq_{a,j}$, so
\[
 3H_p+2S-\sum_jn_jH_j^\Sigma
       =2\sum_{a,j}q_{a,j}d_j-E.
\]
Finally $D_p+\sum_{a,j}q_{a,j}d_j=3T$ and $U_p=2T-\beta$.
Thus $2V\le10T-2\beta-E$, as required. In particular the sign of the
deficit term does not depend on integrality of $n_j$.
\end{proof}

\begin{lemma}[Weighted residual witness budget]\label{lem:residual}
If $k>0$, then for each player $i$
\begin{equation}\label{eq:witness}
                       kH_i\le T-W+(k-q_i)u_i.
\end{equation}
Consequently,
\begin{equation}\label{eq:residualbound}
                         V\le3T+\beta-\frac{C_2}{k}.
\end{equation}
\end{lemma}
\begin{proof}
For every $x\in\core$, fractional team coverage implies
$\sum_ip_ix_i\ge W$. By efficiency,
\begin{equation}\label{eq:residualcover}
                       \sum_iq_ix_i\le T-W=:R.
\end{equation}
In particular $W\le T$. Fix $i\in N_a$ and take a core point with
$x_i=H_i$. Lemma~\ref{lem:range} gives $x_h\ge H_i-u_i$ for every
other $h\in N_a$. Since residual weights and all other-sector payoffs
are nonnegative,
\[
 R\ge\sum_{h\in N}q_hx_h\ge q_iH_i+(k-q_i)(H_i-u_i)
                         =kH_i-(k-q_i)u_i.
\]
Here $k-q_i=\sum_{h\in N_a\setminus\{i\}}q_h\ge0$. The lower bound
$H_i-u_i$ need not be positive, and $q_i$ may be zero.
This proves~\eqref{eq:witness} using just one attaining witness at a time.

Multiply~\eqref{eq:witness} by $q_i$ and sum over players. The resulting
scalar inequalities concern the fixed numbers $H_i$, so different witnesses
can be used. Since $\sum_iq_i=3k$,
\[
                 kS\le3k(T-W)+k\beta-C_2.
\]
Dividing by $k$ proves~\eqref{eq:residualbound}.
\end{proof}

\begin{proof}[Proof of the upper assertion in Theorem~\ref{thm:main}]
For $k>0$, average~\eqref{eq:aggregate} and~\eqref{eq:residualbound}:
\begin{equation}\label{eq:strictbound}
                        V\le4T-\frac E4-\frac{C_2}{2k}.
\end{equation}
If $C_2>0$, this is strictly less than $4T$. If $C_2=0$, each
nonnegative summand $q_i^2u_i$ is zero, hence each $q_iu_i$ is zero
and $\beta=0$. Equation~\eqref{eq:residualbound} then gives
$V\le3T<4T$. If $k=0$, every $q_i=0$ and
\eqref{eq:residualcover} gives $W\le T$, again implying $V\le3T<4T$.

Thus $V<4T$ at every point of the compact polytope $\pack$.
The objective $V$ is linear in $\lambda$ and attains its maximum.
Apply the strict inequality at a maximizing point and use
Lemma~\ref{lem:dual} to conclude $s(M)<T/3$.
\end{proof}

\section{A sharp bound for integral obstructions}\label{sec:integral}
For an integral player-disjoint partial matching $P$, let $U$ be its
unmatched players, $k=n-|P|$, and set
\[
              J(M,P)=v(P)+\frac13\sum_{i\in U}H_i-T.
\]
The indicator of $P$ is feasible in $\pack$, so $J(M,P)\le s(M)$.
An upper bound on $J(M,P)$ alone is not an upper bound on the full subsidy.

\begin{theorem}[Integral residual-count bound]\label{thm:integral}
For $k\ge1$,
\begin{equation}\label{eq:integral}
             J(M,P)\le\frac{2(k-1)}{3(2k-1)}T.
\end{equation}
For $k=0$, $J(M,P)\le0$. For every integer $k\ge2$, equality
in~\eqref{eq:integral} occurs in a market with $n=2k-1$.
\end{theorem}
\begin{proof}
For the indicator packing, $q_i\in\{0,1\}$ and $C_2=\beta$. The two
aggregate bounds imply
\[
               V\le3T+\frac{k-1}{k}\beta,\qquad V\le5T-\beta.
\]
Weight these inequalities by $k/(2k-1)$ and $(k-1)/(2k-1)$,
respectively, and add. This gives
$V\le(4-1/(2k-1))T$ and hence~\eqref{eq:integral}.
The argument includes $k=1$, where the second weight is zero.
For $k=0$, $P$ is perfect and $J(M,P)=v(P)-T\le0$.
Equality for each $k\ge2$ follows from the family below with
$\ell=k-1$.
\end{proof}

\section{The sharpness family}\label{sec:family}
For an integer $\ell\ge1$, take $n=2\ell+1$ and labels
$0,\ldots,2\ell$. Initialize $A,B,C$ to the identity matrix. Define
special teams
\begin{equation}\label{eq:family}
 e_0=(0,\ell,2\ell),\qquad
 e_s=(s,\ell+s,s-1)\quad(1\le s\le\ell-1).
\end{equation}
For every $e=(j_1,j_2,j_3)$ in this list set
$A_{j_1j_2}=2$, $B_{j_1j_3}=1$, and $C_{j_2j_3}=2$.
All labels within a special team are distinct, so these entries are
off-diagonal. The special teams are player-disjoint: their sectorwise
label sets are $\{0,\ldots,\ell-1\}$,
$\{\ell,\ldots,2\ell-1\}$, and
$\{2\ell,0,\ldots,\ell-2\}$. In particular, special entries do not
conflict or overwrite diagonal entries.

\begin{proposition}[Exact values]\label{prop:family}
The market in~\eqref{eq:family} has a nonempty original core and
\begin{equation}\label{eq:familyvalues}
            T=3(2\ell+1),\qquad H_i=3\ (i\in N),\qquad s(M)=2\ell.
\end{equation}
\end{proposition}
\begin{proof}
Consider the bipartite incidence graph with one vertex for each label
and one vertex for each special team, joined by label membership.
Consecutive special teams share one label; nonconsecutive special teams
share none, and each has a private second-sector label. Their label sets
together contain every label $0,\ldots,2\ell$. Thus the graph is connected
and has $(2\ell+1)+\ell$ vertices and $3\ell$ edges, so it
is a tree. This also covers $\ell=1$.

Fix a label $r$ and a sector $a_0\in\{1,2,3\}$. Root this incidence
tree at $r$, assign $r$ color $a_0$, and traverse outward. When a special
team is reached, its parent label is already colored; assign each of its
two other labels the sector in which that label occurs in the team.
No assignments conflict because the incidence graph is a tree. Each
special team therefore has at least two labels colored by their designated
sectors. With $c(j)$ denoting the resulting label color, set
\begin{equation}\label{eq:colorwitness}
           x_{a,j}=\one_{\{a=a_0\}}+2\one_{\{c(j)=a\}}.
\end{equation}
This vector is nonnegative, has total $3n$, pays each diagonal team
exactly $3$, and satisfies $x_{a_0,r}=3$.

We check every full-team inequality. Each team receives a baseline payoff
of $1$ from the first term in~\eqref{eq:colorwitness}. If it contains two
positive off-diagonal pair entries, the shared player lies in the same
special team for both entries, because special teams are player-disjoint.
The team must therefore be that special team. It has value $5$ and receives
at least $1+2+2=5$. If two positive pair entries are diagonal, all three
labels coincide, and the team's value and payoff are both $3$.

Every remaining team has at most one positive off-diagonal entry and at
most one diagonal entry. Any off-diagonal entry belongs to a special team.
At least one of its endpoints has its designated color, since at least two
of the three labels in that special team do. Such a team receives at least
$3$. Its value is at most $3$ if the off-diagonal entry is $2$, and at
most $2$ if that entry is $1$. A nondiagonal team with no positive
off-diagonal entries has value at most $1$, covered by the baseline.
These cases exhaust all teams.

The diagonal matching has value $3n$ and $x$ covers every matching with
total $3n$, so $T=3n$ and $x\in\core$. Diagonal tightness bounds every
original-core coordinate by $3$. As $r,a_0$ were arbitrary,
\eqref{eq:colorwitness} proves $H_i=3$ for each player.

Every protected allocation pays at least $1$ to each player. The $\ell$
special teams form a partial matching of value $5\ell$, and its
$3n-3\ell$ unmatched players require at least $3n-3\ell$. Hence the
total payment is at least $3n+2\ell$, proving $s(M)\ge2\ell$.
Conversely start with payoff $1$ to every player and add $2$ to the
first-sector player of each special team. A special team receives $5$;
every other team has value at most $3$ by the preceding classification
and receives at least $3$. This allocation covers all coalitions, meets
every floor, and has total $3n+2\ell$. Thus $s(M)=2\ell$.
\end{proof}

The ratios $2\ell/[3(2\ell+1)]$ tend to $1/3$, proving the remaining
assertion of Theorem~\ref{thm:main}. For the special partial matching,
$k=\ell+1$ and $J(M,P)=2\ell$, which proves the equality assertion
of Theorem~\ref{thm:integral}.

\begin{remark}[Why the finite bound is integral]\label{rem:fractional}
The formula in~\eqref{eq:integral} cannot be applied by inserting a
fractional packing's residual mass $k$. In the family with $\ell=1$,
let $D$ be the diagonal matching and $P=\{e_0\}$. The packing
$\lambda=\tfrac12\one_D+\tfrac12\one_P$ has $T=9$, $W=7$, and
$k=1$. Its six residual players have $q_i=1/2$ and $H_i=3$, so $S=9$
and $W+S/3-T=1$. Substituting $k=1$ in~\eqref{eq:integral} would
incorrectly give zero. The fractional proof instead retains $C_2$,
which need not equal $\beta$.
\end{remark}

\section{Relation to subsidy frameworks and scope}\label{sec:framework}
Let $l_i=H_i/3$. From diagonal caps, $\sum_i l_i\le T$.
Following Meir et al.~\citep[Section~2.2]{Meir2011} and
Bachrach et al.~\citep[Section~5]{Bachrach2018}, let
$\operatorname{CoS}(x,G)$ be the least subsidy allowing a stable payoff
vector above an efficient preimputation $x$, with $\sum_ix_i=T$.
The following identity is a direct comparison of feasible sets:
\begin{equation}\label{eq:cos}
       s(M)=\min_{x\ge l,\ x(N)=T}\operatorname{CoS}(x,G).
\end{equation}
Indeed, if $y\ge l$ is stable with total at least $T$, there is a vector
$x$ between $l$ and $y$ with total $T$: distribute $T-l(N)$ of the
available nonnegative coordinate increments $y-l$. Conversely, a stable
extension of any such $x$ lies above $l$. In particular the intervention
fits an existing subsidy framework. Its sharp quantitative behavior here
comes from the structure of the assignment market and its original maxima.

Our universal statement allows arbitrary nonnegative real pair values,
dense supports, repeated diagonal labels within a team, and zero diagonal
values. It does not assume that all maxima on one sector are jointly
attainable. Nonempty original core is an explicit hypothesis, not a
claim for all pair-additive three-sided markets. Pair-additivity and
nonnegativity enter through the incident-budget identity
\eqref{eq:budget} and the completion argument in Lemma~\ref{lem:range}.
Equal sector sizes give the perfect diagonal matching and common residual
mass. No extension to signed pair values, arbitrary triple arrays, or a
different productive-coalition convention is asserted.

The finite integral bound does not determine the largest full subsidy at
each fixed $n$, because its maximizing dual may be fractional. A constructive
allocation formula avoiding general linear programming, and the sharp
tradeoff for floors $\alpha H_i$ at other values of $\alpha$, are separate
questions. They are not required for the sharp constant proved here.

\paragraph{Verification and reproducibility.}
The proofs above are conventional all-parameter arguments. The accompanying
standard-library Python package verifies the coefficient identities in
\eqref{eq:coefficient} and the two slack identities in Online Resource~1,
including all $27$ ordered label triples over three symbols in the sixfold
count (covering every coincidence type).
For $\ell=1,\ldots,12$, it reconstructs each separate ideal-attaining
core witness and a matching optimal subsidy allocation from the raw pair
matrices. That check covers $3{,}603{,}804$ full-team inequalities and rejects
$12$ deliberately damaged allocations. Additional rational packing checks
include the fractional counterexample in Remark~\ref{rem:fractional}.
These finite checks corroborate, but do not establish, the universal theorem;
no proof-assistant formalization is claimed.

\paragraph{Supplementary material.}
Online Resource~1 (\texttt{ESM\_1.pdf}) gives the slack identities, a notation
recap, and the scope of the exact tests. Online Resource~2
(\texttt{ESM\_2.zip}) contains the standalone verification programs, frozen
inputs, source hashes, and the fresh execution report. The full October~7
compendium remains a separate research record.

\section*{Statements and Declarations}
\paragraph{Author information, funding and interests.}
Author names, order, affiliation or residence, corresponding contact, funding
and competing interests require confirmation for this paper before submission.
No declaration of no funding or no competing interests is presumed.
\paragraph{Data and code availability.}
No empirical dataset is used. Exact mathematical inputs and verification code
are included in Online Resource~2. This preparation has not newly deposited
or submitted the manuscript or its materials.
\paragraph{AI assistance and responsibility.}
OpenAI Codex and Anthropic Claude assisted exploration, proof drafting,
code and internal review. Internal model review and finite checks are not
human peer review or formal proof verification. The confirmed author(s)
must assess the arguments and references, confirm declarations, and assume
responsibility before submission. Historical priority and human specialist
assessment remain open.

\begin{thebibliography}{99}
\small
\bibitem[Atay and N\'u\~nez(2019)]{AtayNunez2019}
A. Atay and M. N\'u\~nez.
Multi-sided assignment games on $m$-partite graphs.
\emph{Annals of Operations Research} 279:271--290, 2019.
\href{https://doi.org/10.1007/s10479-019-03256-5}{https://doi.org/10.1007/s10479-019-03256-5}.

\bibitem[Bachrach et al.(2018)]{Bachrach2018}
Y. Bachrach, E. Elkind, E. Malizia, R. Meir, D. Pasechnik,
J.~S. Rosenschein, J. Rothe, and M. Zuckerman.
Bounds on the cost of stabilizing a cooperative game.
\emph{Journal of Artificial Intelligence Research} 63:987--1023, 2018.
\href{https://doi.org/10.1613/jair.1.11270}{https://doi.org/10.1613/jair.1.11270}.

\bibitem[Lucas(1995)]{Lucas1995}
W.~F. Lucas. Core theory for multiple-sided assignment games.
\emph{Duke Mathematical Journal} 81(1):55--65, 1995.
\href{https://doi.org/10.1215/S0012-7094-95-08106-X}{https://doi.org/10.1215/S0012-7094-95-08106-X}.

\bibitem[Mart\'inez-de-Alb\'eniz et al.(2026)]{Martinez2026}
F.~J. Mart\'inez-de-Alb\'eniz, C. Rafels, and N. Ybern.
Maximum and minimum core payoffs for supermodular multisided assignment
markets with local monotonicity.
\emph{Discrete Applied Mathematics} 379:605--612, 2026.
\href{https://doi.org/10.1016/j.dam.2025.10.006}{https://doi.org/10.1016/j.dam.2025.10.006}.

\bibitem[Meir et al.(2011)]{Meir2011}
R. Meir, J.~S. Rosenschein, and E. Malizia.
Subsidies, stability, and restricted cooperation in coalitional games.
\emph{Proceedings of IJCAI 2011}, pp.~301--306.
\href{https://www.cs.huji.ac.il/~jeff/papers/coopmas11meir.pdf}{Author-hosted extended manuscript, Section 2.2}.

\bibitem[N\'u\~nez Lugilde et al.(2024)]{Priority2024}
I. N\'u\~nez Lugilde, A. Est\'evez-Fern\'andez, and E. S\'anchez-Rodr\'iguez.
Priority coalitional games and claims problems.
\emph{Mathematical Methods of Operations Research} 100:669--701, 2024.
\href{https://doi.org/10.1007/s00186-024-00876-x}{https://doi.org/10.1007/s00186-024-00876-x}.
The source comparison uses the accepted manuscript dated May 22, 2024.

\bibitem[Quint(1991)]{Quint1991}
T. Quint. The core of an $m$-sided assignment game.
\emph{Games and Economic Behavior} 3(4):487--503, 1991.
\href{https://doi.org/10.1016/0899-8256(91)90017-9}{https://doi.org/10.1016/0899-8256(91)90017-9}.

\bibitem[Shapley and Shubik(1971)]{ShapleyShubik1971}
L.~S. Shapley and M. Shubik.
The assignment game I: The core.
\emph{International Journal of Game Theory} 1:111--130, 1971.
\href{https://doi.org/10.1007/BF01753437}{https://doi.org/10.1007/BF01753437}.

\bibitem[Stuart(1997)]{Stuart1997}
H.~W. Stuart, Jr. The supplier--firm--buyer game and its $m$-sided generalization.
\emph{Mathematical Social Sciences} 34(1):21--27, 1997.
\href{https://doi.org/10.1016/S0165-4896(96)00830-X}{https://doi.org/10.1016/S0165-4896(96)00830-X}.
\end{thebibliography}
\end{document}
